% =====================================================================
% SECTION I -- INTRODUCTION (style mirrors the companion E2E-CDRO paper)
% =====================================================================
\section{Introduction}
\label{sec:introduction}

\subsection{Motivation and Challenges}
\label{subsec:motivation}

Artificial intelligence (AI) data centers are becoming large, fast-growing,
and increasingly grid-relevant electricity loads, and power systems have begun
to demand flexibility in return for interconnection. Texas Senate Bill~6
subjects large loads to curtailment-readiness requirements~\cite{sb6_2025},
hyperscalers have signed demand-response agreements covering on the order of
a gigawatt of data-center load~\cite{google_dr2025}, and modest flexibility
from large new loads would unlock tens of gigawatts of interconnection
capacity in existing U.S. systems~\cite{norris2025rethinking}. The supply side of this bargain is
forming on the \emph{compute} side: a field demonstration in Phoenix reduced a
production AI cluster's power by $25\%$ for three hours by orchestrating
workloads~\cite{colangelo2025phoenix}, and workload-level studies report
$18$--$55\%$ power flexibility across AI job types~\cite{acun2026flexibility}.

This paper develops the complementary resource: the \emph{cooling side} of
the liquid-cooled AI facility, field-demonstrated for over a decade but
never committed as a certified market product. Direct-to-chip liquid cooling
concentrates a substantial, controllable share of facility power in a
chilled-water plant coupled to megajoules of coolant and facility-water
thermal mass. Measurements on a production liquid-cooled hall show that
relaxing the coolant supply temperature from $17$ to $25^\circ$C sheds $63\%$
of cooling power~\cite{gheni2026liquid}; field tests a decade ago already
exercised cooling-side demand response in air-cooled
facilities~\cite{ghatikar2012demand}, and co-adjusting chilled-water setpoints
with server power yields measurable regulation capacity~\cite{fu2020assessments}.
Unlike batteries or workload curtailment, this thermal inertia is already
paid for, and it responds in minutes.

What is missing is the step from \emph{potential} to \emph{product}. A
day-ahead market participant must commit a deliverable quantity hours before
delivery, yet the cooling plant's deliverable depth is not a number but a set.
Answering how an operator can commit this set safely leads to two key
challenges.

\textit{1) The deliverable depth is context-dependent and weather-coupled.}
The anti-condensation limit places a floor on the coolant supply temperature
just above the ambient dew point, so humidity directly truncates the
admissible cooling reductions; workload volatility consumes the same thermal
buffer the product sells; and the depth that can be \emph{sustained} decays
with the product duration. A deterministic maximum commitment invites
delivery failures and hotspot excursions, while a static worst-case
commitment forfeits most of the value. For variable generation and
distributed energy resources (DERs), the analogous problem is resolved by
the do-not-exceed (DNE) limit~\cite{zhao2015variable} and the dynamic
operating envelope~\cite{aemo2022doe}, and the buildings literature
certifies reserves against a static comfort band; none covers a feasible
set whose lower boundary is an exogenous stochastic weather variable.

\textit{2) A day-ahead commitment must survive real-time uncertainty with
unequal stakes.} Between commitment and delivery, the realized workload and
dew point deviate from their forecasts, and the two consequences differ in
kind: a delivery shortfall is a \emph{priced} event settled by penalties,
whereas a hotspot or condensation excursion is a hardware-safety event that no
payment compensates. A single robustness level cannot serve both: tightening
until safety holds in the worst case destroys the offer, while sizing
uncertainty for the average case lets safety erode exactly when the product
is exercised.

\begin{table*}[t]
\centering
\caption{Positioning of ENCORE against representative flexibility
characterization and commitment methods}
\label{tab:positioning}
\renewcommand{\arraystretch}{1.08}
\setlength{\tabcolsep}{2.2pt}
\scriptsize
\begin{tabularx}{\textwidth}{
>{\hsize=1.45\hsize\linewidth=\hsize\raggedright\arraybackslash}X
>{\hsize=0.85\hsize\linewidth=\hsize\centering\arraybackslash}X
>{\hsize=1.05\hsize\linewidth=\hsize\raggedright\arraybackslash}X
>{\hsize=0.80\hsize\linewidth=\hsize\centering\arraybackslash}X
>{\hsize=1.00\hsize\linewidth=\hsize\raggedright\arraybackslash}X
>{\hsize=0.85\hsize\linewidth=\hsize\centering\arraybackslash}X
}
\toprule
\textbf{Method category}
& \textbf{Resource}
& \textbf{Flexibility object}
& \textbf{Weather coupling}
& \textbf{Uncertainty treatment}
& \textbf{\mbox{Delivery guarantee}} \\
\midrule
Data-center DR field studies~\cite{ghatikar2012demand,fu2020assessments}
& Cooling (air)
& Measured potential
& No
& None
& No \\
Compute-side flexibility~\cite{colangelo2025phoenix,acun2026flexibility,wierman2014opportunities}
& IT workload
& Demonstrated reduction
& No
& Operational
& No \\
Thermal storage retrofits~\cite{oh2025rtes,zhang2025carnot,ravi2026prosumer}
& Added hardware
& Storage dispatch
& Partial
& Scenario-based
& No \\
DC bidding / workload commitment~\cite{chen2023bilevel,fan2026commitment}
& IT workload
& Day-ahead bid
& No
& Stochastic program
& Partial \\
Buildings-to-grid reserves~\cite{vrettos2016scheduling,vrettos2016robust,gorecki2017experimental}
& HVAC mass
& Reserve capacity
& Comfort band (static)
& Robust MPC
& Yes (reserve) \\
DNE limits / operating envelopes~\cite{zhao2015variable,aemo2022doe}
& Generation / DER
& Dispatch envelope
& No
& Robust set
& Yes (dispatch) \\
TCL virtual battery / flexibility function~\cite{hao2015aggregate,junker2018characterizing,warrington2013policy}
& TCL aggregate
& Battery equivalent
& No
& Aggregate bounds
& Partial \\
\textbf{ENCORE (this work)}
& \textbf{Liquid cooling loop}
& \textbf{Polyhedral deliverability envelope}
& \textbf{Dew-point-coupled}
& \textbf{Nested conformal sets}
& \textbf{Two-clause certificate} \\
\bottomrule
\end{tabularx}
\vspace{-5mm}
\end{table*}

\vspace{-4mm}
\subsection{Related Work}
\label{subsec:related}

\textit{1) Data centers as grid resources.} Surveys and system studies
established schedulable IT load as a flexibility
source~\cite{takci2025datacentres,wierman2014opportunities}, and recent work
demonstrates it at production scale on AI
clusters~\cite{colangelo2025phoenix,acun2026flexibility}; bi-level bidding and
day-ahead workload commitment with deliverability
constraints~\cite{chen2023bilevel,fan2026commitment} bring the compute side
into market form. All act on the IT workload. On the cooling side, early
field studies exercised demand response in air-cooled
facilities~\cite{ghatikar2012demand}, frequency-regulation capacity was
obtained by co-adjusting chilled-water setpoints with server
power~\cite{fu2020assessments}, and recent measurements quantify large, fast
supply-temperature flexibility in production liquid-cooled
halls~\cite{gheni2026liquid}. These works quantify \emph{potential}; none
defines a committable product with a delivery guarantee. Thermal-storage
retrofits~\cite{oh2025rtes,zhang2025carnot,ravi2026prosumer} monetize heat by
adding hardware, whereas here the unmodified loop itself is the resource.

\textit{2) Buildings-to-grid reserve provision.} The methodological template
for offering thermal flexibility under uncertainty comes from commercial
buildings: robust capacity allocation ahead of time, an online model
predictive control (MPC) layer, and experimental market
demonstrations~\cite{vrettos2016scheduling,vrettos2016robust,gorecki2017experimental}.
ENCORE descends from this line and departs where the resource does: the
feasible set's lower boundary is an exogenous stochastic variable (the dew
point), not a static comfort band; the dynamics are minutes-scale
loop physics against seconds-scale GPU bursts, not hours-scale
building mass; and the deliverable is the certified envelope object itself,
not a reserve schedule.

\textit{3) Certifying the maximum offer.} Do-not-exceed limits for variable
resources~\cite{zhao2015variable}, dynamic operating envelopes for
distribution-connected resources~\cite{aemo2022doe}, aggregate virtual
batteries for thermostatic loads~\cite{hao2015aggregate}, flexibility
functions~\cite{junker2018characterizing}, and policy-based
reserves~\cite{warrington2013policy} all certify an admissible operating
region, not a point forecast. ENCORE is the thermal-resource member of
this family. Within it, the combination of an \emph{exact} deliverability
polyhedron, a context-conditional (conformal~\cite{lei2018distribution})
uncertainty description, and a settlement that prices the residual delivery
risk is, to our knowledge, new. Relative to the companion contextual robust
\emph{control} work~\cite{shen2026cdro}, which optimizes operating cost
through an online min--sup layer, the decision object here is a market
commitment with a certificate: robustness acts offline as constraint
tightening of an envelope, and the market layer (product, settlement,
penalty) is part of the model. Table~\ref{tab:positioning} summarizes the
positioning.

\vspace{-3mm}
\subsection{Our Contributions}
\label{subsec:contributions}

ENCORE lets a liquid-cooled AI data center sell its cooling-side thermal
inertia in a day-ahead market with guarantees that survive to real-time
delivery. The main contributions are as follows.

\begin{itemize}
\item We formulate cooling-side flexibility as a day-ahead offering problem
with an explicit product, settlement, and penalty structure. The formulation
couples the (offer, duration) product to hotspot safety, an anti-condensation
supply-temperature floor, actuator limits, and whole-hour energy metering,
under joint workload and dew-point uncertainty.

\item We characterize the deliverable (offer, duration) set \emph{exactly} as
a context-dependent polyhedron, the certified flexibility envelope, and
derive a closed-form weather-coupled virtual-battery equivalent whose
capacity, power limits, and leakage are functions of the dew-point forecast.
We further show that consecutive full-depth delivery is infeasible under
whole-hour energy settlement, so certified commitments must be spaced.

\item We develop a two-stage offering method with end-to-end guarantees.
Nested context-conditional disturbance sets tighten the envelope offline
through exact tube margins: a safety set, never traded for revenue, and a
delivery set whose level is a product design variable priced by the
settlement. A certified fallback policy transfers both guarantees to the
real-time layer. We prove that any
offer inside the tightened envelope satisfies thermal safety on the safety
set and meets its settled obligation at confidence $1-\epsd$, with expected
penalties bounded.

\item We validate the certificate end-to-end on measured records: a
2.0-million-worker GPU-cluster trace with a causal job-aware day-ahead
forecast, archived numerical-weather-prediction (NWP) dew-point forecasts,
and ten 2024 weeks of prices from the Electric Reliability Council of Texas
(ERCOT). Across ten held-out weeks $\times$ $20$ activation draws the
certified participant incurs zero in-domain violation episodes (against
$229$ violation day-seeds for the uncertified envelope), delivery failures
stay inside the priced budget, and day-ahead workload information alone
raises the certifiable offer by $83\%$.
\end{itemize}

Section~\ref{sec:problem_formulation} models the hall and the offering
problem; Section~\ref{sec:methodology} develops ENCORE and its guarantees;
Section~\ref{sec:experimental_evaluation} evaluates;
Section~\ref{sec:conclusion} concludes.
